\section{Languages and bridges}
\label{sec:languages}

A \emph{language} in this library is a syntactic class of game
descriptions paired with a compiler $\compileTo : \mathrm{Syntax} \to
\KG$ that fixes its denotational semantics.  Equipping the syntax
with an SOS-style operational semantics is optional but
ergonomic:\ several theorems about a language are most naturally
stated by induction on the SOS, and SOS-equivalence between source
programs is then a structural certificate that they compile to the
same kernel game.

The library exposes five languages — normal-form (NFG),
extensive-form (EFG), multi-agent influence diagrams (MAID),
multi-round games (MR, with stochastic and repeated specializations),
and factored-observation stochastic games (FOSG) — and the
relation-parameterized expressiveness framework that relates them
(\S\ref{sec:languages:expressiveness}).

\subsection{Normal-form games}
\label{sec:languages:nfg}

The normal-form layer is the simplest:\ a normal-form game over a
finite player type $\Players$ and a finite action family
$\Act_i$ is the data
\[
  \NFG = \bigl(\,\Players,\,\{\Act_i\}_i,\,\util : ({\textstyle\prod_i}\Act_i)
                  \to \Payoff{\Players}\,\bigr),
\]
i.e., a function from joint pure actions to payoff vectors.  The
compilation $\NFG \compileTo \KG$ takes the strategy type to be
$\Act_i$, the outcome type to be $\prod_i \Act_i$, the kernel to
be the Dirac evaluation $\Profile \mapsto \Dirac_\Profile$, and the
utility to be the input function.  The kernel-game image is
deterministic; randomness in NFGs enters only through the mixed
extension (\S\ref{sec:semantic-core:kernelgame}).

The library's NFG layer also formalizes two specialized
strategic-form structures whose form is a refinement of NFG and
whose theorems are about that refinement:\ \emph{congestion
games}~\citep{Rosenthal1973}, where a strategy is a subset of
shared resources and a player's payoff sums per-resource delays
that depend only on congestion (and Rosenthal's potential makes
every congestion game an exact potential game, hence yields pure
Nash existence and the finite improvement property); and
\emph{public goods}, two free-standing theorems on the
parameterized payoff $e + \mathrm{mpcr} \cdot \sum_j c_j - c_i$
(free-riding is dominant when $\mathrm{mpcr} < 1$, full
cooperation Pareto-dominates full defection when $\mathrm{mpcr}
\cdot n > 1$).  A \emph{Stackelberg leader-follower} structure
sits in the same directory, although strictly speaking it is
sequential rather than strategic-form:\ the leader commits, the
follower best-responds, and the leader-advantage theorem
(Stackelberg payoff $\ge$ simultaneous Nash payoff) is proved
against the follower-BR function.

Beyond these specialized structures, the NFG layer carries a
catalogue of canonical 2-player worked examples:\ Prisoner's
Dilemma, Matching Pennies, Stag Hunt, Hawk–Dove (Chicken), Battle
of the Sexes, Dictator Game, Traveler's Dilemma (two-action
fragment), Cournot duopoly with three quantities,
Bertrand price competition with three discrete prices (which has
both a strict undercut-resistant Nash and a weak price-floor Nash
under discrete pricing), and a two-game Braess paradox instance
(\TT{braessRestricted}/\TT{braessAugmented} showing welfare
strictly decreases when a dominant strategy is added).  Each
defines an action type with \TT{deriving Fintype, DecidableEq}, a
payoff table by pattern match, and one or two pure-Nash facts
discharged by case analysis (\TT{cases i <;> cases a' <;> simp;
try linarith}); Matching Pennies additionally has an exact mixed
characterization in \TT{Languages/NFG/MatchingPenniesMixed.lean}.
That file instantiates the language-independent
\TT{Concepts/BinaryMixed.lean} theorem for binary
matching-pennies-like \KGs, showing that both players put probability
$1/2$ on heads. These serve both as reader-facing examples and as
smoke tests for the NFG$\to$\KG\ compilation pipeline. The
Stackelberg leader-follower file additionally carries an
entry-deterrence instance (\TT{EntryDeterrence.game}) where the
leader's fight-commitment strictly dominates their
accommodate-and-Nash payoff of 1, illustrating the abstract
leader-advantage theorem on a concrete payoff matrix.

The cooperative-flavored game-theoretic models that used to live
in \TT{Languages/NFG/} — TU coalitional games with the Shapley
value~\citep{Shapley1953}, Nash bargaining
\citep{Nash1950bargaining}, and stable matching predicates and
structural lemmas~\citep{GaleShapley1962} — are not strategic-form
formalisms (they take coalitions, feasible payoff sets, or
preference rankings as primitive rather than per-player
strategies) and have been moved to a separate
\TT{Cooperative/} branch outside the \texttt{Languages/NFG/}
hierarchy.  See \S\ref{sec:future:cooperative} for what they
contain.

\leancorr{\TT{Languages/NFG/} contains \TT{Syntax}, \TT{Compile},
\TT{Examples} (canonical 2-player games), \TT{Tests},
\TT{Theorems}, plus the strategic-form subclass files
\TT{CongestionGame}, \TT{PublicGoods}, \TT{Stackelberg}.
The cooperative models live under \TT{Cooperative/} as
\TT{CoalitionalGame}, \TT{Bargaining}, \TT{Matching}.
Supplement~§4.1.}

\subsection{Extensive-form games}
\label{sec:languages:efg}

An \emph{extensive-form game} is a finite rooted tree whose
non-terminal nodes are owned either by Nature (chance nodes with
a probability distribution over children) or by a player (decision
nodes from which the player chooses one child).  Terminal nodes
carry payoffs.  Two pieces of textbook structure on top of this
tree are crucial.  First, the assignment of decision nodes to
\emph{information sets}:\ a partition of each player's decision
nodes into equivalence classes the player cannot tell apart at
play time.  When a player moves at a node in an information set
$I$, they know they are at some node of $I$, but not which one;
their strategy must therefore commit to the same action across
all nodes of $I$.  This is what makes EFGs an
imperfect-information formalism — perfect-information games are
the special case where every information set is a singleton.

\begin{center}
\begin{tikzpicture}[
  font=\small,
  node distance=8mm,
  inner/.style={circle, draw, fill=black!8, minimum size=4mm, inner sep=1pt},
  term/.style={rectangle, draw=none, font=\footnotesize},
  edgelbl/.style={font=\scriptsize, midway},
]
\node[inner] (p1)              {1};
\node[inner] (p2L) at (-2,-1)  {2};
\node[inner] (p2R) at ( 2,-1)  {2};
\node[term]  (o1)  at (-2.8,-2.1) {$(1,0)$};
\node[term]  (o2)  at (-1.2,-2.1) {$(0,1)$};
\node[term]  (o3)  at ( 1.2,-2.1) {$(0,1)$};
\node[term]  (o4)  at ( 2.8,-2.1) {$(1,0)$};
\draw (p1) -- node[edgelbl, above left]{$L$} (p2L);
\draw (p1) -- node[edgelbl, above right]{$R$} (p2R);
\draw (p2L) -- node[edgelbl, left]{$\ell$} (o1);
\draw (p2L) -- node[edgelbl, right]{$r$} (o2);
\draw (p2R) -- node[edgelbl, left]{$\ell$} (o3);
\draw (p2R) -- node[edgelbl, right]{$r$} (o4);
\draw[dashed, rounded corners=2pt]
  ($(p2L)+(-0.4,0.4)$) rectangle ($(p2R)+(0.4,-0.4)$);
\node[font=\scriptsize] at (0,-0.45) {information set};
\end{tikzpicture}
\end{center}

\noindent
Player~1 moves first ($L$ or $R$);  player~2 moves second but
without observing 1's choice — so 2's two decision nodes are in
one information set (the dashed box), and 2's pure strategies
are just $\{\ell, r\}$, not the four functions $\{L,R\}\to\{\ell,r\}$.

The library's EFG follows the textbook structure closely but
factors it slightly differently for the type theory:\ rather than
parameterize the tree by a player type directly, players are
factored through an \emph{information structure}:
\begin{equation}
\label{eq:infostructure}
  \InfoStruct = \bigl(\,n,\,\InfoSet : \mathrm{Fin}\,n \to \mathsf{Type},\,
                       \mathrm{arity} : \forall p\,I,\,\NN_{>0}\,\bigr),
\end{equation}
mapping each player $p \in \mathrm{Fin}\,n$ to a finite type of
information sets and each information set to a positive action
arity.  A \emph{game tree} over $\InfoStruct$ is the inductive
type
\begin{equation}
\label{eq:gametree}
\begin{aligned}
  \Tree(\InfoStruct, \Outcome) \;\coloneqq\;\;
    &\mathsf{terminal}(z : \Outcome) \\
  \mid\;& \mathsf{chance}(k, \mu : \PMF{\mathrm{Fin}\,k}, \mathrm{next} : \mathrm{Fin}\,k \to \Tree) \\
  \mid\;& \mathsf{decision}(p, I : \InfoSet_p, \mathrm{next} : \mathrm{Fin}\,(\mathrm{arity}_p\,I) \to \Tree).
\end{aligned}
\end{equation}
A decision node stores only the information-set ID; the player
identity is recovered from the structure.  An EFG is a tree paired
with a utility $\Outcome \to \Payoff{\Players}$.  Strategies come
in three flavors:
\begin{itemize}\itemsep2pt
\item \emph{Pure}: $\Spure_p = \forall I : \InfoSet_p,\, \mathrm{Fin}\,(\mathrm{arity}_p\,I)$.
\item \emph{Behavioral}: $\Sbeh_p = \forall I,\, \PMF{\mathrm{Fin}\,(\mathrm{arity}_p\,I)}$.
\item \emph{Mixed}: $\Smix_p = \PMF{\Spure_p}$.
\end{itemize}

The library's EFG semantics is \emph{distributional}:\ given a
behavioral profile $\beta : \forall p,\, \Sbeh_p$, the function
$\evalDist(\beta, t) \in \PMF{\Outcome}$ is defined by structural
recursion on $t$, with chance nodes binding the chance distribution
and decision nodes binding the player's behavioral distribution at
the relevant information set.  The compilation
$\EFG \compileTo \KG$ takes strategies to be behavioral profiles,
outcomes to be the EFG's $\Outcome$, the kernel to be $\evalDist$,
and the utility to be the EFG utility.

\designnote{Distributional, not deterministic, semantics.  An
alternative would be to evaluate to a deterministic outcome under
a pure profile and lift to mixed/behavioral profiles externally.
We chose distributional semantics because it (a) fits chance nodes
without ad-hoc lifting, (b) makes Kuhn's theorem
(\S\ref{sec:theorems:kuhn}) a purely structural identity
$\evalDist(\beta) = \evalDist(\mathrm{realize}(\mu))$ for the
realization $\mu$ of $\beta$, and (c) fits the
$\K : \prod_i \Strat_i \to \PMF{\Outcome}$ shape of \KG\ without
the strategies-to-outcomes adapter.}

An \emph{augmented} EFG in the FOSG-paper sense — a thin wrapper
over an EFG that adds a public partition over all reachable
histories and a per-player information partition over reachable
histories, including those at which the player does not act — is
also provided.  Standard EFG information sets are only defined at
decision nodes for the acting player; the augmentation lifts this
to all of history space and is the carrier on which the FOSG/EFG
bridge (\S\ref{sec:languages:bridges}) is most naturally stated.

The EFG \TT{Examples} file carries two canonical sequential games:
the 2-stage \emph{Centipede game} (Rosenthal 1981; payoffs (1,0)
/ (0,2) / (3,1) with backward-induction prediction `(p0 takes)`
leaving the welfare-max (3,1) on the table), and the
\emph{Ultimatum game} (Güth--Schmittberger--Schwarze 1982; a
proposer offers one of three splits to a responder who
accepts/rejects). Each ships with the four `evalPure` outcome
identities plus the textbook backward-induction comparison
inequalities; the full
\TT{EFGGame.IsSubgamePerfectEq} witnesses are not packaged.

\leancorr{\TT{Languages/EFG/} contains \TT{Syntax}, \TT{SOS},
\TT{Compile}, \TT{CompileObs} (compilation to ObsModelCore for
Kuhn), \TT{Augmented} (the augmented EFG), \TT{Kuhn},
\TT{Refinements} (subgame perfection,
\TT{EFGGame.IsSubgamePerfectEq}), \TT{Examples} (binary choice
plus the centipede and ultimatum), \TT{Tests}.
Supplement~§4.2.}

\subsection{Multi-agent influence diagrams}
\label{sec:languages:maid}

A \emph{multi-agent influence diagram} (MAID,
\citealp{KollerMilch2003}) represents a game as a directed acyclic
graph whose nodes are typed as one of three kinds:
\begin{itemize}\itemsep2pt
\item \emph{Chance nodes} $X$ — random variables; each carries a
  conditional probability distribution (CPD) over its domain given
  the realizations of its parents.  This is Nature's contribution.
\item \emph{Decision nodes} $D$ — each owned by a specific player;
  the parents of a decision node represent the information the
  owning player observes when deciding, and the player chooses a
  value in the node's finite domain as a function of those
  parents.
\item \emph{Utility nodes} $U$ — real-valued functions of their
  parents; each utility node is associated with a specific player,
  and that player's payoff is the sum of all their utility nodes.
\end{itemize}

\noindent
The DAG structure makes causal and informational dependence
\emph{explicit} in the syntax:\ the parent set of a decision node
is literally the information available to the deciding player, and
the parent set of a utility node is literally the variables it
depends on.  Where an EFG has to flatten the same information into
a tree with shared information sets, a MAID exposes the same
content as a Bayesian network with decision and utility annotations.

\begin{center}
\begin{tikzpicture}[
  font=\small,
  node distance=12mm,
  chance/.style ={circle,   draw, fill=black!5,  minimum size=8mm, inner sep=1pt},
  decision/.style={rectangle,draw, fill=black!10, minimum size=8mm, inner sep=1pt},
  utility/.style ={diamond,  draw, fill=black!8,  minimum size=8mm, inner sep=1pt},
  arr/.style={-{Stealth[length=2mm]}, thick, black!75},
]
\node[chance]   (M)  at (0, 1.4)   {$M$};
\node[decision] (D1) at (-2,0)     {$D_1$};
\node[decision] (D2) at ( 2,0)     {$D_2$};
\node[utility]  (U1) at (-2,-1.4)  {$U_1$};
\node[utility]  (U2) at ( 2,-1.4)  {$U_2$};
\draw[arr] (M)  -- (D1);
\draw[arr] (M)  -- (D2);
\draw[arr] (D1) -- (U1);
\draw[arr] (D1) -- (U2);
\draw[arr] (D2) -- (U1);
\draw[arr] (D2) -- (U2);
\draw[arr] (M)  to[bend right=40] (U1);
\draw[arr] (M)  to[bend  left=40] (U2);
\node[font=\scriptsize, anchor=west] at (3.0, 1.4)   {chance};
\node[font=\scriptsize, anchor=west] at (3.0, 0)     {decision (player 1, 2)};
\node[font=\scriptsize, anchor=west] at (3.0,-1.4)   {utility (player 1, 2)};
\end{tikzpicture}
\end{center}

\noindent
A two-player MAID:\ a market state $M$ is sampled by Nature; each
player observes $M$ and chooses an action $D_i$; each player's
payoff $U_i$ depends on the market state and both decisions.
Decision nodes are squares, chance nodes are circles, utility
nodes are diamonds.

The library formalizes a finite-domain fragment:\ each node has a
finite domain, decision-node parents specify the player's
information, and utility nodes are real-valued.

The compilation $\MAID \compileTo \KG$ proceeds via an
intermediate \emph{frontier evaluator}.  At each step, the
evaluator picks a node whose parents are all assigned and either
(a) samples its chance distribution, (b) consults the player's
strategy at the appropriate information set, or (c) accumulates
its utility into the per-player running sum.  The order of
evaluation is non-trivial:\ when several decision nodes share a
parent set, the evaluator's resolution must be order-independent
to make the compilation well-defined.  The diamond/commutation
properties licensing the order-independent compiler are proved
once; the same \emph{frontier-stability} pattern reappears in the
event-graph-based protocol layer that depends on this library
\citep{KollerMilch2003}.

\tradeoff{The order-of-evaluation issue is, in our view, the
single largest design pressure in the MAID layer.  An obvious
alternative — fixing a topological order at compilation time —
would simplify the proofs but build into the formalism a choice
that no MAID-as-a-mathematical-object cares about.  We accepted
the proof obligations and got an order-independent compiler in
exchange.}

\leancorr{\TT{Languages/MAID/} contains \TT{Syntax}, \TT{SOS},
\TT{FrontierEval}, \TT{FrontierLemmas} (the diamond properties),
\TT{FoldEval}, \TT{Prefix}, \TT{Compile}, \TT{CompileObs},
\TT{Kuhn}, \TT{Theorems}.  Supplement~§4.3.}

\subsection{Multi-round games, repeated games, stochastic games}
\label{sec:languages:multiround}

A multi-round game is a finite list of
\emph{rounds}~\citep{ShohamLeytonBrown2008}, each round consisting
of a state-dependent signal distribution (nature's joint signal),
a per-player view function (state plus signal mapped to an
observation), and a deterministic state transition keyed on the
joint action profile.  Strategies are pure and memoryless at the
protocol level — a function from the player's view to an
\emph{optional} action — with mixed/behavioral and recall
structure imposed externally as analysis-layer concerns.

The library's compilation $\MR \compileTo \KG$ generates the
outcome kernel by folding the round list and threading the state
through the chain of observations and transitions.  A bounded
horizon ensures termination.  Around this layer the library also
provides two standard variants:
\begin{itemize}\itemsep2pt
\item \emph{Repeated games}:\ identical rounds with public
observation of past action profiles.  \TT{RepeatedGame.lean}
provides an infinite-horizon discounted repeated-game helper:
strategies map finite public histories to the next stage action,
payoffs are normalized discounted sums, and
\TT{constStrategy\_isDiscountedNash} records the forward direction
that a stage-game Nash action profile remains Nash under constant
repetition, assuming bounded stage payoffs.

The infinite discounted repeated-game folk theorem is not built
through the finite-list protocol syntax; it is formalized directly
on public \TT{KernelGame} profile streams in
\TT{Concepts/FolkTheorem.lean}.  In the main theorem the repeated
game is taken over the mixed extension, so period actions are
publicly observed mixed stage strategies.  In that observable
mixed-action model, every strictly individually rational feasible
payoff vector is approximated by Nash payoffs of sufficiently
patient discounted repeated games (\Cref{thm:folk-discounted}).

\item \emph{Stochastic games}~\citep{Shapley1953stochastic}:\
rounds chosen by state, with a stochastic transition.  Bounded
horizon, by design.
\end{itemize}

The multi-round layer also compiles into the generic
\emph{InfoModel}/\emph{ObsModel} substrate that the
behavioral$\,\leftrightarrow\,$mixed Kuhn theorem
(\S\ref{sec:theorems:kuhn}) is stated against.  This is the layer
at which a per-round full-recall hypothesis becomes a per-round
\emph{ActionPosteriorLocal} hypothesis sufficient for the
mixed-to-behavioral direction.

\leancorr{\TT{Languages/MultiRound/} contains \TT{Syntax},
\TT{SOS}, \TT{Compile}, \TT{CompileObs} and its linearizations
(\TT{CompileObsLin}, \TT{CompileObsLinAdequacy},
\TT{CompileObsLinRecall}), \TT{RepeatedGame}, \TT{StochasticGame},
\TT{Kuhn}, \TT{Theorems}.  Supplement~§4.4.}

\subsection{Factored-observation stochastic games}
\label{sec:languages:fosg}

To explain factored-observation stochastic games we have to
distinguish two pre-existing notions.  A \emph{stochastic
game}~\citep{Shapley1953stochastic} is a Markov game on a finite
state space:\ at each state, every active player simultaneously
chooses an action; the joint action determines a stage payoff and
a stochastic transition to a next state, and play continues to
horizon.  This formalism handles state and transition randomness
cleanly but is opaque about \emph{what each player observes}
during a transition — the textbook treatment assumes either
perfect monitoring of the state or a particular monitoring scheme
named in prose.  A \emph{factored-observation stochastic game}
(FOSG, \citealp{Kovarik2022}) makes the monitoring explicit and
\emph{factored}:\ each transition produces a per-player
\emph{private} observation (independent component per player) and
a single \emph{public} observation visible to everyone.  A player's
strategy is then a function of their own observation \emph{sequence}
— the only data the player has access to.

\begin{center}\small
\begin{tabular}{@{}lp{0.62\textwidth}@{}}
\toprule
Component & What it adds \\
\midrule
state space + transition & ``where am I, and where am I going?'' \\
active-set per state & which players must move at each state \\
per-state available actions & legality of each player's action \\
reward vector & per-step payoff to each player \\
per-player private obs & what player $i$ alone sees on a transition \\
public observation & what \emph{every} player sees on a transition \\
\bottomrule
\end{tabular}
\end{center}

\noindent
The FOSG framework subsumes the standard formalisms used in the
multi-agent reinforcement-learning literature (POMGs,
Dec-POMDPs) by an appropriate choice of private and public
observation channels, and it specializes to extensive-form games
when the per-player observation sequence is rich enough to identify
the current information set.  This is the modern carrier for
imperfect-information sequential games:\ a player's strategy can
depend only on their own observation sequence, which is
\emph{exactly} the formal counterpart of an information set in the
sequential setting.

In the library, an FOSG carries a state space, an active-set
function, per-player available actions, a transition kernel, a
reward vector, per-player private observations, and a public
observation.  Histories are recorded by an explicit history type.
The finite-horizon FOSG kernel game is the canonical native
carrier on which the bounded Kuhn theorem
(\S\ref{sec:theorems:kuhn}) is proved.

\leancorr{\TT{Languages/FOSG/} contains \TT{Basic}, \TT{History},
\TT{Information}, \TT{Strategy}, \TT{Values}, \TT{Execution},
\TT{OutcomeClosure}, \TT{Serial}, \TT{Compile}, \TT{Kuhn},
\TT{Theorems}.  Supplement~§4.5.}

\subsection{Intrinsic-form games}
\label{sec:languages:intrinsic}

A focused layer formalizes the \emph{intrinsic} (a.k.a.\
Witsenhausen-style) representation of games following
\citet{HeymannDeLaraChancelier2020}:\ rather than fixing an order
in which agents act, the syntax exposes a set of agent decision
\emph{addresses} and information sets directly.  Nature first
selects an exogenous state; choices at the addresses then have to
fit together into a global configuration.  A model is
\emph{solvable} when, for every state and every profile of local
choices, there is a unique compatible configuration and hence a
well-defined law over realized decisions.

This syntax supports the same three ways of randomizing that
appear in sequential games, but without building a tree.  A
player may draw one complete contingent plan, randomize
independently across that player's decision addresses, or randomize
independently at the information atoms the player can distinguish.  The
first comparison is unconditional:\ independent randomization by
address and behavioral randomization by information atom induce
the same local choice laws.  The second is the intrinsic analogue
of Kuhn's theorem.  If the information structure admits a causal
order and the player has perfect recall along that order, then
the conditional behavioral strategy extracted from a mixed plan
preserves the full outcome law against every opponent profile.
The conclusion is equality of distributions over realized
decision profiles, not just equality of one-address marginals.

The intrinsic layer is a separate syntactic root than
EFG/MAID/MR/FOSG, used as a target by representations that prefer
order-independent reasoning.

\leancorr{\TT{Languages/Intrinsic/} contains the strategy
carriers, causal-order and perfect-recall predicates, and the
formal outcome-law equivalence for solvable W-models.  The
definition and theorem anchors are listed in Supplement~§4.6.}

\subsection{Bridges as commuting compilations}
\label{sec:languages:bridges}

A \emph{bridge} between two languages $L_1, L_2$ is a translation
$\tau : \mathrm{Syntax}_{L_1} \to \mathrm{Syntax}_{L_2}$ paired with
a theorem stating that the two compilation paths agree under a
specified semantic relation:
\begin{equation}
\label{eq:bridge}
  R\bigl(\,L_1.\compileTo(x),\;L_2.\compileTo(\tau(x))\,\bigr).
\end{equation}
Different relations $R$ correspond to different bridge strengths:
EU-isomorphism (\S\ref{sec:semantic-core:morphisms}) is the
strongest; protocol simulation (utility-preserving but not
necessarily strategy-bijective) is intermediate;
EU-equivalence-of-equilibria is the minimum useful target.  The
library exposes the following bridges:

\begin{center}\small
\begin{tabular}{@{}llp{0.42\textwidth}@{}}
\toprule
Source & Target & Relation \\
\midrule
NFG  & EFG  & EU-isomorphism (NFG as a depth-1 EFG) \\
NFG  & FOSG & EU-isomorphism (NFG as a one-shot FOSG) \\
EFG  & NFG  & EU-equivalence on the realized strategic form \\
MAID & EFG  & strategic-form refinement (information sets preserved) \\
EFG  & FOSG & game-form simulation (FOSG view of an EFG with serial execution) \\
\bottomrule
\end{tabular}
\end{center}

The bridge theorems are stated as identities on the compiled
kernel games (or game forms), and are the glue that lets a
theorem about, say, EFG behavioral strategies be transported to a
MAID by composing the MAID$\to$EFG bridge with the relevant
strategic-form theorem on the EFG side.

\leancorr{\TT{Languages/Bridges/} contains \TT{NFG\_EFG},
\TT{NFG\_FOSG}, \TT{EFG\_NFG}, \TT{MAID\_EFG}, and \TT{FOSG/}
(\TT{ObsModelCore}, \TT{AugmentedEFG}, \TT{SerialExec}).
Supplement~§5.2.}

\subsection{Relation-parameterized expressiveness}
\label{sec:languages:expressiveness}

The library packages the bridges into a generic expressiveness
framework parameterized by a semantic relation $R$ on
$\KG$.\footnote{In the Lean code the relation, together with its
\TT{refl}/\TT{symm}/\TT{trans} proofs, is bundled in a record called
\TT{EquivalenceLens}; we mention the name only because the supplement
uses it and so does any reader who reads the source.  ``Lens'' here
is the project's bundling vocabulary; it carries no connection to
optics in functional programming or to lenses in category theory.}
A \emph{language} in this layer is a pair $(\mathrm{Syntax},
\compileTo)$ with $\compileTo$ targeting $\KG_\Players$ for some
fixed $\Players$; a \emph{reduction} $L \to M$ over a relation $R$
is a translation $\tau$ and a proof that
$R(L.\compileTo(x), M.\compileTo(\tau(x)))$ holds for every source
$x$.  Reductions compose under transitivity of $R$.

\paragraph{The relation taxonomy.}\
The library catalogues the semantic relations this framework
supports.  They form a refinement order:\ a stronger relation
determines a weaker one.

\begin{center}\small
\begin{tabular}{@{}lp{0.62\textwidth}@{}}
\toprule
Relation & What it preserves \\
\midrule
$\mathrm{ProtocolEquivalent}$ &
  Protocol isomorphism on game forms (strategies and outcomes
  relabeled bijectively, kernel preserved); ignores utility. \\
$\mathrm{UtilityDistributionEquivalent}$ &
  Bisimulation in \KG: both protocol and joint utility
  distribution preserved at every profile. \\
$\mathrm{UtilityDistributionSimulation}$ &
  One-way version of the above; the directed reduction. \\
$\mathrm{EUEquivalent}$ &
  Equal expected utility for each player at corresponding profiles
  (the equilibrium-relevant relation). \\
$\mathrm{EUSimulation}$ &
  Directed EU-comparison; equilibria of the source pull back to
  the target. \\
\bottomrule
\end{tabular}
\end{center}

\noindent
The bridges of \S\ref{sec:languages:bridges} are reductions in
this framework for the relation in their column.  Context-aware
extensions of the framework support languages whose compilation
depends on a typing/visibility context — the natural target for
downstream protocol-language formalizations.

\paragraph{Worked theorems.}\
Two concrete worked expressiveness theorems are packaged in this
form.  An expressiveness claim like ``every finite NFG is
a one-shot FOSG up to EU-isomorphism'' becomes the statement that
$\NFG \to \FOSG$ is a reduction over the EU-isomorphism relation;
``every bounded MAID is an EFG up to strategic-form refinement''
becomes the analogous reduction over the refinement relation.

\designnote{Expressiveness comparisons in the protocol-language
literature are notoriously informal.  By making the relation $R$
explicit, the library forces every comparison statement to commit
to which structural information is preserved and which is not.
The resulting lattice of bridges (NFG, EFG, MAID, MR, FOSG ordered
by which relation we can prove between them) is itself a
contribution.}

\leancorr{\TT{Languages/Expressiveness/} contains \TT{Basic}
(\TT{Language}, \TT{Reduction}), \TT{Relations} (the relation
taxonomy), \TT{CausalContext}, \TT{MAID\_EFG}, \TT{EFG\_FOSG}.
Supplement~§5.1.}
